\documentclass[12pt]{article}
\usepackage[utf8]{inputenc}
\usepackage[T1]{fontenc}
\usepackage{hyperref}
\title{Lachirioag Zapotec Language Resources for K-12 Curriculum Development}
\author{Compiled by Claude Sonnet 5.5 (claude-sonnet-5-5) for Andr\'as Kornai}
\date{}
\begin{document}
\maketitle

\section{Introduction}
The code ztc (Glottocode lach1251) is retired. Glottolog records that Lachirioag Zapotec was merged into Yatee Zapotec (zty), effective 2007-07-18, and treats it as a dialect of Yatee. Yatee Zapotec is spoken in Oaxaca, Mexico, and Wikipedia gives about 5,000 speakers as of 2004. No resources specific to Lachirioag were found. Any materials would fall under Yatee Zapotec (zty), which would need a separate search.

\section{Primary Educational Resources}
None found for Lachirioag.

\section{Linguistic and Reference Materials}
\begin{itemize}
\item Glottolog entry for Lachirioag Zapotec, recording the retirement and merger into zty. \url{https://glottolog.org/resource/languoid/id/lach1251}
\end{itemize}

\section{Religious and Community Resources}
None found for Lachirioag.

\section{Digital and Audio Resources}
\begin{itemize}
\item Global Recordings Network page for Yatee Zapotec (zty), which lists Lachirioag Zapotec as an alternate name. \url{https://www.globalrecordings.net/en/language/zty}
\end{itemize}

\section{Curriculum Development Context}
Because ztc is a retired code, work on this variety should be done under Yatee Zapotec (zty). Community-specific materials, for example from San Crist\'obal Lachiruaj, would need to be located with local partners.

\end{document}